% Section 2: the data structure, the generation algorithm, and the proof that
% what comes out is the spanning tree Section 1 said it had to be. The cost of
% all of this is deferred to Section 6; only correctness is argued here.
\section{Generating the maze into a graph}
\label{sec:maze-generation}

Section~\ref{sec:problem-formulation} fixed the target: a spanning tree of
$\gridgraph$, drawn at random. This section builds one. The structure comes
first, because the algorithm is written against it, and the guarantee comes
last, because it is the construction rule that supplies it.

\subsection{What has to be stored}

Two different things are worth keeping, and confusing them is the usual way to
end up storing neither well.

\begin{description}[leftmargin=*,style=nextline]
  \item[The maze itself --- which doors exist.] A subset of $\gridedges$. It is
    what a walker needs, and the only thing needed to answer Problem~2.
  \item[The way it was built --- the region decomposition.] Which wall was put
    where, and in what order. It is the binary tree the divide-and-conquer rule
    produces for free, and the object whose depth the cost analysis of
    Section~\ref{sec:analysis} is written in terms of.
\end{description}

\subsection{The cell array: two links per cell}

The doors are held in a $W\times H$ array $T$ whose entry $T[x][y]$ is one cell
record holding two links,
\[
  T[x][y]=\bigl\langle P_0,\;P_1 \bigr\rangle,
\]
where $P_0$ points at the cell $(x+1,y)$ to its right and $P_1$ at the cell
$(x,y+1)$ below it, and where a link set to \textsc{nil} means \emph{no door
there}. Figure~\ref{fig:cell-record} draws the arrangement.

% Why two links per cell are enough: each cell owns the edge to its right
% neighbour and the edge to the one below, so every edge of the grid graph is
% stored exactly once and the boundary cells are where the null links live.
\begin{figure}[htbp]
  \centering
  \begin{tikzpicture}[scale=1.0]
    \foreach \x in {0,1,2} \foreach \y in {0,1,2} {
      \node[cell] (c\x\y) at (1.7*\x, -1.7*\y) {\scriptsize$(\x,\y)$};
    }
    % P0: the link to the right neighbour
    \foreach \x/\xn in {0/1, 1/2} \foreach \y in {0,1,2}
      \draw[link] (c\x\y) -- node[above, font=\tiny, doorgreen] {$P_0$} (c\xn\y);
    % P1: the link to the neighbour below
    \foreach \y/\yn in {0/1, 1/2} \foreach \x in {0,1,2}
      \draw[link] (c\x\y) -- node[right, font=\tiny, doorgreen] {$P_1$} (c\x\yn);
    % the boundary cells own nothing in the direction that leaves the area
    \foreach \y in {0,1,2}
      \draw[link, wallred, dash pattern=on 1.5pt off 1.5pt]
        (c2\y) -- ++(0.95,0) node[right, font=\tiny] {\textsc{nil}};
    \foreach \x in {0,1,2}
      \draw[link, wallred, dash pattern=on 1.5pt off 1.5pt]
        (c\x2) -- ++(0,-0.95) node[below, font=\tiny] {\textsc{nil}};
  \end{tikzpicture}
  \caption{The cell record. Every cell stores two links, $P_0$ to the cell on
  its right and $P_1$ to the cell below it; the cells of the last column and
  the last row carry \textsc{nil} in the direction that leaves the area. Each
  edge of $\gridgraph$ is owned by exactly one cell, so the $2WH-W-H$ edges are
  held without a single one stored twice, and the four-neighbour view a search
  wants is read back in $O(1)$ from the left and upper neighbours' own links.}
  \label{fig:cell-record}
\end{figure}


Two links, not four, is the one design decision in the structure. Every edge of
$\gridgraph$ joins a cell to the cell on its right or the cell below it, never
to a third kind of neighbour, so \emph{ownership by the upper-left endpoint}
covers the whole edge set exactly once: the $2WH-W-H$ possible doors are held
in $2WH$ slots with no edge duplicated, and the last column and last row are
where the unused slots sit.

\begin{remark}
The four-neighbour record --- a node holding its coordinates and an array
$\mathit{Neighbour}[4]$ --- is the convenient \emph{view} for a search, which
wants all ways out of a cell at once. It is not a second structure: the
neighbours of $(x,y)$ to the west and north are recovered in $O(1)$ from
$T[x-1][y].P_0$ and $T[x][y-1].P_1$. Materialising it costs twice the memory
and buys no information, so this report stores two links and reads four.
\end{remark}

\subsection{The region tree}

The second structure is a binary tree, one node per call of the recursion:

\begin{itemize}[leftmargin=*]
  \item an \textbf{internal node} records one cut --- its orientation $t$, the
        boundary $j$ the wall sits on, the position $i$ of its single door ---
        and has the two regions the cut created as its children;
  \item a \textbf{leaf} is a region reduced to one cell.
\end{itemize}

Figure~\ref{fig:region-tree} shows the first three cuts of a $4\times4$ area
beside the top of the tree they build. Since every cut creates exactly one
door, internal nodes and doors are in bijection: a tree with $WH$ leaves has
$WH-1$ internal nodes, which is the $|E|=WH-1$ of
Proposition~\ref{prop:spanning-tree} arrived at from the other side.

% The picture and the data structure are the same object seen twice: every wall
% implied on the left is one internal node on the right, and the parameters
% labelling that node are exactly the t, j and i the algorithm drew.
\begin{figure}[htbp]
  \centering
  \begin{subfigure}[b]{0.36\textwidth}\centering
    \begin{tikzpicture}[scale=0.78]
      \fill[cellgrey] (0,0) rectangle (4,-4);
      \draw[black!60] (0,0) rectangle (4,-4);
      \draw[black!20] (1,0) -- (1,-4); \draw[black!20] (3,0) -- (3,-4);
      \draw[black!20] (0,-1) -- (4,-1); \draw[black!20] (0,-2) -- (4,-2);
      \draw[black!20] (0,-3) -- (4,-3);
      % division A: the root, a vertical wall at j=1, door at row i=1
      \draw[wall segment] (2,0) -- (2,-1);
      \draw[wall segment] (2,-2) -- (2,-4);
      \draw[door] (2,-1) -- (2,-2);
      \node[wallred, font=\scriptsize\bfseries] at (2.35,0.28) {A};
      % division B: inside the left half, horizontal wall at j=1, door at column 0
      \draw[wall segment] (1,-2) -- (2,-2);
      \draw[door] (0,-2) -- (1,-2);
      \node[wallred, font=\scriptsize\bfseries] at (-0.32,-2) {B};
      % division C: inside the right half, horizontal wall at j=2, door at column 3
      \draw[wall segment] (2,-3) -- (3,-3);
      \draw[door] (3,-3) -- (4,-3);
      \node[wallred, font=\scriptsize\bfseries] at (4.32,-3) {C};
    \end{tikzpicture}
    \caption{The first three divisions of a $4\times4$ area.}
    \label{fig:region-tree-cuts}
  \end{subfigure}\hfill
  \begin{subfigure}[b]{0.60\textwidth}\centering
    \begin{tikzpicture}[scale=1.0, level distance=13mm,
      level 1/.style={sibling distance=42mm},
      level 2/.style={sibling distance=21mm}]
      \node[tree node] {A: $t{=}1,\,j{=}1,\,i{=}1$}
        child {node[tree node] {B: $t{=}0,\,j{=}1,\,i{=}0$}
          child {node[tree leaf] {$\region{0,0}{1,1}$}}
          child {node[tree leaf] {$\region{0,2}{1,3}$}}
        }
        child {node[tree node] {C: $t{=}0,\,j{=}2,\,i{=}3$}
          child {node[tree leaf] {$\region{2,0}{3,2}$}}
          child {node[tree leaf] {$\region{2,3}{3,3}$}}
        };
    \end{tikzpicture}
    \caption{The top of the region tree it builds. The four boxes at the bottom
    are regions still to be cut, not cells.}
    \label{fig:region-tree-tree}
  \end{subfigure}
  \caption{Every division is one internal node. The node records the
  orientation $t$, the wall position $j$ and the door position $i$ it drew, and
  its two children are the two regions it created; a leaf is a corridor. A tree
  with $L$ leaves has $L-1$ internal nodes, so of the $WH-1$ doors of
  Proposition~\ref{prop:spanning-tree}, $L-1$ are division doors and the rest
  are laid down inside the corridors.}
  \label{fig:region-tree}
\end{figure}


\subsection{Starting full, then deleting}

The algorithm starts from $T$ with \emph{every} link set: that is $\gridgraph$
itself, all doors open, every cell reachable from every other. Generation only
ever deletes. This is the direction that makes the invariant cheap to state and
to keep:

\begin{quote}
  \emph{At every call, the region handled is connected in $T$, and no edge of
  $T$ joins that region to anything outside it that the caller did not put
  there.}
\end{quote}

It holds at the root because $\gridgraph$ is connected, and
Figure~\ref{fig:one-partition-step} is one step of preserving it.

\subsection{One partition step}

A call handles a region $\region{a_0,a_1}{b_0,b_1}$. If the region is a single
cell there is nothing to cut and the call returns a leaf. Otherwise:

\begin{enumerate}[leftmargin=*]
  \item \textbf{Choose the orientation $t$.} The handout asks for a wall across
        the longest dimension, so $t=1$ (a vertical wall, cutting the column
        range) when $a_1-a_0>b_1-b_0$, $t=0$ (a horizontal wall, cutting the
        row range) when it is smaller, and a random bit decides a tie. The rule
        never picks a dimension of extent $0$: a region one cell wide is cut
        across the other dimension, and a region one cell wide in both is the
        base case.
  \item \textbf{Choose the wall position $j$}, uniformly among the interior
        boundaries of the chosen dimension, and \textbf{the door position $i$},
        uniformly along the wall.
  \item \textbf{Delete every crossing but one.} Walk $k$ along the wall and
        delete the edge it crosses at $k$, skipping $k=i$.
  \item \textbf{Recurse} on the two regions, and hang them under the node
        recording $\langle t,j,i\rangle$.
\end{enumerate}

% One turn of the loop in Algorithm 2, shown on the graph rather than on the
% drawing: the region arrives with every crossing edge present and leaves with
% exactly one of them kept.
\begin{figure}[htbp]
  \centering
  \begin{subfigure}[b]{0.47\textwidth}\centering
    \begin{tikzpicture}[scale=0.82]
      \foreach \x in {0,1,2,3,4} \foreach \y in {0,1,2,3}
        \node[node vertex] (u\x\y) at (\x,-\y) {};
      \foreach \x/\xn in {0/1,1/2,2/3,3/4} \foreach \y in {0,1,2,3}
        \draw[black!45] (u\x\y) -- (u\xn\y);
      \foreach \y/\yn in {0/1,1/2,2/3} \foreach \x in {0,1,2,3,4}
        \draw[black!45] (u\x\y) -- (u\x\yn);
    \end{tikzpicture}
    \caption{On entry the region is the full grid graph on
    $\region{a_0,a_1}{b_0,b_1}$: connected, and full of cycles.}
    \label{fig:partition-before}
  \end{subfigure}\hfill
  \begin{subfigure}[b]{0.47\textwidth}\centering
    \begin{tikzpicture}[scale=0.82]
      \foreach \x in {0,1,2,3,4} \foreach \y in {0,1,2,3}
        \node[node vertex] (w\x\y) at (\x,-\y) {};
      \foreach \x/\xn in {0/1,1/2,3/4} \foreach \y in {0,1,2,3}
        \draw[black!45] (w\x\y) -- (w\xn\y);
      \foreach \y/\yn in {0/1,1/2,2/3} \foreach \x in {0,1,2,3,4}
        \draw[black!45] (w\x\y) -- (w\x\yn);
      \draw[door] (w21) -- (w31);
      \draw[wall] (2.5,0.45) -- (2.5,-0.75);
      \draw[wall] (2.5,-1.25) -- (2.5,-3.45);
      \node[font=\tiny, wallred, above] at (2.5,0.45) {wall at $j=2$};
      \node[font=\tiny, doorgreen, left] at (-0.2,-1) {door at $i{=}1$};
    \end{tikzpicture}
    \caption{On exit every edge crossing column boundary $j$ is gone except the
    one at row $i$.}
    \label{fig:partition-after}
  \end{subfigure}
  \caption{One partition step, with $t=1$ (a vertical wall). The loop walks
  $k$ over the rows $b_0..b_1$ of the region and deletes the edge
  $(j,k)\!-\!(j{+}1,k)$ for every $k\neq i$. The two halves stay connected
  internally, and the single surviving edge is the only way between them ---
  which is exactly the induction step of
  Proposition~\ref{prop:generation-is-a-tree}.}
  \label{fig:one-partition-step}
\end{figure}


\subsection{The algorithm}

\begin{algorithm}[t]
\caption{Maze generation by recursive division.}
\label{alg:generate-maze}
\SetKwFunction{GenerateMaze}{GenerateMaze}
\SetKwFunction{CarveRegion}{CarveRegion}
\SetKwFunction{SplitOrientation}{SplitOrientation}
\SetKwFunction{FullGrid}{FullGridGraph}
\SetKwFunction{RandomIn}{RandomIn}
\SetKwFunction{RandomBit}{RandomBit}
\SetKwFunction{Delete}{DeleteDoor}
\SetKwFunction{Leaf}{Leaf}
\SetKwFunction{NodeOf}{CutNode}
\KwIn{$W,H\ge1$, the width and height of the area in cells.}
\KwOut{$T$, the cell array holding a spanning tree of $\gridgraph$; $r$, the
       root of the region tree that built it.}
\BlankLine
\Fn{\GenerateMaze{$W$, $H$}}{
  $T \leftarrow$ \FullGrid{$W$, $H$}\tcp*{all doors open}
  $r \leftarrow$ \CarveRegion{$T$, $[0,W-1]$, $[0,H-1]$}\;
  \Return $(T,r)$\;
}
\BlankLine
\Fn{\CarveRegion{$T$, $[a_0,a_1]$, $[b_0,b_1]$}}{
  \lIf{$a_0=a_1$ \textup{\textbf{and}} $b_0=b_1$}{\Return \Leaf{$a_0$, $b_0$}}
  $t \leftarrow$ \SplitOrientation{$a_1-a_0$, $b_1-b_0$}\;
  \eIf{$t=1$}{
    $j \leftarrow$ \RandomIn{$a_0$, $a_1-1$}\tcp*{wall between columns $j$ and $j+1$}
    $i \leftarrow$ \RandomIn{$b_0$, $b_1$}\tcp*{the one door, on row $i$}
    \For{$k \leftarrow b_0$ \KwTo $b_1$}{
      \lIf{$k=i$}{\textbf{continue}}
      \Delete{$T$, $(j,k)$, $(j+1,k)$}\;
    }
    $\ell \leftarrow$ \CarveRegion{$T$, $[a_0,j]$, $[b_0,b_1]$}\;
    $r \leftarrow$ \CarveRegion{$T$, $[j+1,a_1]$, $[b_0,b_1]$}\;
  }{
    $j \leftarrow$ \RandomIn{$b_0$, $b_1-1$}\tcp*{wall between rows $j$ and $j+1$}
    $i \leftarrow$ \RandomIn{$a_0$, $a_1$}\tcp*{the one door, on column $i$}
    \For{$k \leftarrow a_0$ \KwTo $a_1$}{
      \lIf{$k=i$}{\textbf{continue}}
      \Delete{$T$, $(k,j)$, $(k,j+1)$}\;
    }
    $\ell \leftarrow$ \CarveRegion{$T$, $[a_0,a_1]$, $[b_0,j]$}\;
    $r \leftarrow$ \CarveRegion{$T$, $[a_0,a_1]$, $[j+1,b_1]$}\;
  }
  \Return \NodeOf{$t$, $j$, $i$, $\ell$, $r$}\;
}
\BlankLine
\Fn{\SplitOrientation{$d_x$, $d_y$}}{
  \lIf{$d_x>d_y$}{\Return $1$}
  \lIf{$d_x<d_y$}{\Return $0$}
  \Return \RandomBit{}\;
}
\end{algorithm}

\Delete{} clears one link: deleting the edge $(j,k)\!-\!(j{+}1,k)$ sets
$T[j][k].P_0 \leftarrow \textsc{nil}$, and deleting $(k,j)\!-\!(k,j{+}1)$ sets
$T[k][j].P_1 \leftarrow \textsc{nil}$. Ownership by the upper-left endpoint
means each deletion touches one slot and no lookup is needed to find it.

The exit plays no part in generation. Proposition~\ref{prop:spanning-tree}
makes every cell reachable from every other, so any boundary cell may be
declared $\exitcell$ afterwards --- which is also why the same maze serves the
many trap-door entries the situation asks for.

\subsection{Why the result is a spanning tree}

\begin{proposition}[Correctness of Algorithm~\ref{alg:generate-maze}]
\label{prop:generation-is-a-tree}
On a region of $m$ cells, \textup{\texttt{CarveRegion}} leaves the subgraph of
$T$ induced on that region connected and acyclic, with exactly $m-1$ doors. In
particular \textup{\texttt{GenerateMaze}} returns a spanning tree of
$\gridgraph$.
\end{proposition}

\begin{proof}
By induction on $m$. If $m=1$ the region is one cell: no edges, connected and
acyclic, $m-1=0$ doors, and the call returns without touching $T$.

Let $m>1$. The orientation rule picks a dimension of extent at least $1$, so
$j$ and the two sub-regions exist and are non-empty; call them $R_\ell$ and
$R_r$, with $m_\ell+m_r=m$ cells. Every edge of the region that crosses the
wall does so at exactly one position $k$ along it, and the loop deletes all of
them but the one at $k=i$. No edge internal to $R_\ell$ or to $R_r$ is touched,
so the two recursive calls receive their regions with every internal door still
open and, by the induction hypothesis, return them connected, acyclic and
carrying $m_\ell-1$ and $m_r-1$ doors.

The only remaining edge between the two is the door at $i$. Joining two
disjoint trees by a single edge yields a tree: it is connected, since any two
cells are joined inside their own half or by crossing that one door, and
acyclic, since a cycle would have to cross the wall an even, hence non-zero,
number of times while only one crossing exists. The door count is
$(m_\ell-1)+(m_r-1)+1=m-1$.

Applying this to the root region gives a connected acyclic spanning subgraph of
$\gridgraph$ on all $WH$ cells: a spanning tree, with $WH-1$ doors.
\end{proof}

The guarantee is therefore structural. Nothing is checked after the fact, no
attempt is ever rejected and retried, and the uniqueness of the route out ---
the whole of Problem~2 --- is a consequence of the shape of the recursion
rather than of a property that had to be tested.

\subsection{The graph that comes out}

Three readings of the one result, none of them a separate copy:

\begin{center}
\begin{tabular}{@{}lll@{}}
  \toprule
  View & Held as & Used by \\
  \midrule
  Doors, compactly   & the array $T$ of $\langle P_0,P_1\rangle$ records
                     & generation, and storage \\
  Adjacency          & the four-neighbour view read off $T$
                     & the search of Section~\ref{sec:shortest-path} \\
  How it was built   & the region tree rooted at $r$
                     & the cost analysis of Section~\ref{sec:analysis} \\
  \bottomrule
\end{tabular}
\end{center}

Figure~\ref{fig:maze-as-graph} is the output of one $4\times4$ run, whose first
three cuts are the ones drawn in Figure~\ref{fig:region-tree-cuts}: $15$ doors
on $16$ cells, one path between any two of them. What
Section~\ref{sec:unit-edge-weights} adds is a length on those doors, so that
``the path out'' becomes a quantity a search can minimise.
